\documentclass{article}
\usepackage{graphicx} % Required for inserting images
\usepackage{amsmath}
\usepackage[margin=1in]{geometry}
\title{Electrical Circuit Formula}
\author{For preparation formula}
\date{}

\begin{document}

\maketitle

\section{DC circuit Analysis}
\textbf{Current(I)}: Rate of change flow. $$I = \frac{dq}{dt}$$\\
\textbf{voltage(V)}: Energy per unit charge. $V = \frac{dw}{dq}$, where dw = the amount of work done.\\
\textbf{Power(P)}: Rate of energy absorption/supply. $$P = VI = I^2R = \frac{V^2}{R}$$\\
\textbf{Energy(W)}: Integrated power over time. $$W = \int{Pdt}$$
\textbf{Ohm's Law}: $$V = IR$$
\textbf{Kirchhoff's Current Law(KCL)} and \textbf{(KVL)}
\[
\sum I_{in} = \sum I_{out}
\]
\[
\sum V_{loop} = 0
\]
\textbf{ Series and parallel circuits}:\\
Series Resistance: $$R_{eq} = R_1 + R_2 +\dots + R_n $$\\
Parallel Resistance:
\begin{align}
\frac{1}{R_{eq}} &= \frac{1}{R_1}+\frac{1}{R_2}+\cdots+\frac{1}{R_n}\\
          R_{eq} &= \frac{R_1R_2}{R_1 + R_2}
\end{align}
Voltage Divider ( for $R_{1} $ in series with $ R_2$):
\begin{center}
    $V_1 = V_{total}*(\frac{R_1}{R_1 + R_2}) $
\end{center}
Current Divider (for $R_1 parallel to R_2$):
\begin{center}
    $ I_1 = I_{total}*(\frac{R_2}{R_1 + R_2}$
\end{center}
\textbf{Delta-Wye($\Delta$--Y)Conversions}\\
Delta to Wye:
\begin{align}
    R_1 = \frac{R_bR_c}{R_a + R_b + R_c}\\
    R_2 =\frac{R_aR_c}{R_a + R_b + R_c}\\
    R_3 =\frac{R_aR_b}{R_a + R_b + R_c}
\end{align}
Wye to Delta:
\begin{align}
    R_a =\frac{R_1R_2 + R_2R3 + R_3R_1}{R_1}\\
    R_b =\frac{R_1R_2 + R_2R3 + R_3R_1}{R_2}\\
    R_c =\frac{R_1R_2 + R_2R3 + R_3R_1}{R_3}
\end{align}
\textbf{Circuit Network Theorems}:
\begin{itemize}
    \item Source Transformation: $ V_s = I_s*R$ and $ I_s = \frac{V_s}{R}$
    \item Thevenin Voltage and Resistance: $ V_{Th} = V_{OC} $ ; (oc = Open circuit)   
     \begin{center}
        $ R_{Th} = \frac{V_{oc}}{I_{sc}}$ ; (Short-circuit current $ I_{sc}$
     \end{center}
     \item Norton Current And Resistance: $$ I_N = I_{sc}, R_N = R_{Th}$$
\end{itemize}
\textbf{Inductors and Capacitors}
Capacitor(C):
\begin{itemize}
    \item Current: i(t) = C$\frac{dv(t)}{dt}$
    \item Voltage: V(t) = $ \frac{1}{C}\int_{t_0}^{t}i(\tau)d\tau + v(t_0)$
    \item Stored Energy: Wc = $ \frac{1}{2}CV^2$
    \item DC steady State: Acts as an open circuit.
\end{itemize}
Inductor(L):
\begin{itemize}
    \item Voltage: v(t) = L$ \frac{di(t)}{dt}$
    \item Current: i(t) = $ \frac{1}{L}\int_{t_0}^t d\tau + i(t_0)$
    \item Stored Energy: $ W_L = \frac{1}{2}LI^2$
    \item DC Steady State: Acts as a short circuit.
\end{itemize}
* First-order Transient Response (RC and RL Circuits):
\begin{center}
    $x(t) =x(\infty)+ [x(t_0)- x(\infty)]e^{\frac{-t}{\tau} }$
\end{center}
For RC circuits: Time constant $\tau = R\cdot C$
For RL circuits: Time constant $\tau = \frac{L}{R}$
\section{AC Sinusoidal Steady-State Analysis}
\begin{enumerate}
    \item Sinusoidal Basics and Waveform factors
    \begin{itemize}
        \item Sinusoidal Expression: $ v(t) = V_msin(wt + \phi$
        \item Angular Frequency: $ \omega = 2\pi f = \frac{2\pi}{T}$
        \item Root-Mean-Square(RMS): $ V_{rms} = \frac{V_m}{\sqrt{2}}\approx0.707V_m$
        \item Average Value(Half-Cycle for sine wave): $ V_{avg} = \frac{2V_m}{\pi}\approx0.637V_m$
        \item Form Factor: Form Factor = $ \frac{V_rms}{V_{avg}}=\frac{\pi}{2\sqrt{2}}$
        \item Peak Factor: Peak Factor = $\frac{V_m}{V_rms} = \sqrt{2} \approx1.414$
    \end{itemize}
    \item Phasors, Impedance(Z), and Admittance(Y)
    \begin{itemize}
        \item Phasor Transform: $ V(t) = V_mcos(\omega t + \theta) \ and\ V = V_m\angle\theta = V_me^j\theta$
        \item Impedance(Z = R + jX):
        \begin{itemize}
            \item Resistor(R): $ Z_R = R$
            \item Inductor(L): $ Z_L = j\omega L = jX_L \ (Inductive Reactance X_L = \omega L)$
            \item Capacitor(C): $ Z_C = \frac{1}{j\omega C} = -jX_C \\ Capacitance Reactance, X_C = \frac{1}{\omega C}$
        \end{itemize}
        \item Admittance $(Y = G + jB = \frac{1}{Z})$:
        \begin{itemize}
            \item G. Conductance, B.Susceptance.
            \begin{center}
                $ Y = \frac{1}{R + jX} = \frac{R}{R^2 + X^2}-j\frac{X}{R^2 + X^2}$
            \end{center}
        \end{itemize}
            
        \end{itemize}
    \item AC Circuit Combination(R-L, R-C, R-L-C)
        \begin{itemize}
            \item Series RLC Impedance:
            \begin{center}
                $ Z = R + j(\omega L- \frac{1}{\omega C}) = |Z|\angle \theta $ 
                $|Z|=\sqrt{R^2 +(X_L - X_C)^2}, \ \theta = tan^{-1}(\frac{X_L - X_C}{R}) $
            \end{center}
            \item AC Ohm's Law: V = I.Z
        \end{itemize}
    \item Sinusoidal Steady-State Power
    \begin{itemize}
        \item Instantaneous Power: p(t) = v(t)*i(t)
        \item Complex Power(S): $ S = V_{rms}I_{rms} = P + jQ = |S|\angle \theta $
        \item Real/Average Power(P): Measured in Watts(W). $$P = V_{rms}I_{rms}cos\theta = I_{rms}^2R $$
        \item Reactive Power(Q): Measured in Volt-Amperes Reactive(VAR)
        \begin{center}
            $Q = V_{rms}I_{rms}sin\theta = I_{rms}^2X $
        \end{center} 
        \item Apparent Power(S): Measured in Volt-Ampere(VA). $$ S = |S| = V_{rms}I_{rms} = \sqrt{P^2 + Q^2}$$
        \item Power Factor(PF): $ PF = cos\theta = \frac{P}{S}$
    \end{itemize}
    \item AC Maximum Power Transfer and Impedance Matching
    \begin{itemize}
        \item For a load impedance $ Z_L = Z_{Th} \Rightarrow R_L = R_{Th} \ and X_L = -X_{Th} $
        \item Maximum Transferred Real Power: $$ P_{max} \frac{|V_{V_{Th},rms}|}{4R_{Th}} $$
        
    \end{itemize}
\section{Technique of General Circuit Analysis}
\begin{itemize}
    \item \textbf{Advanced Network Analysis Equations}\\
    Node-Voltage Method(KCL Matrix):\\
    When solving a circuit with n nodes using Nodal Analysis, you set a system of linear equations using admittance $(G = \frac{1}{R}) $
    $$
    \begin{bmatrix}
        G_{11} & G_{12} & \dots \\
        G_{21} & G_{22} & \dots
    \end{bmatrix}
    \begin{bmatrix}
        V_1\\
        V_2
    \end{bmatrix} =
    \begin{bmatrix}
        I_1\\
        I_2
    \end{bmatrix}
    $$
    \item Mesh-Current Method(KVL Matrix)\\
    When solving a circuit with m meshes using Mesh Analysis, you set up a system using resistance(R):\\
    $$
    \begin{bmatrix}
        R_{11} & R_{12} & \dots \\
        R_{21} & R_{22} & \dots
    \end{bmatrix}
    \begin{bmatrix}
        I_1\\
        I_2
    \end{bmatrix} =
    \begin{bmatrix}
        V_1\\
        V_2
    \end{bmatrix}
    $$
    \item Superposition Theorem\\
    For any linear variable x(voltage or current) affected by independent sources $ S_1, S_2, \dots , S_n:$
    \begin{center}
        $ x = x_1 + x_2 + \dots + x_n$
    \end{center}
    (Where $x_i$ is the value computed while keeping source $S_i$ active and killing all other independent sources: short-circuiting voltage sources and open-circuiting current sources).
\end{itemize}
\section{Component Combinations (Inductors and Capacitors)}
\begin{itemize}
    \item Inductors(L) 
    \begin{center}
        $ Series: L_{eq} = L_1 + L_2 + \dots + L_n $ \\
        $ Parallel: \frac{1}{L_{eq}} = \frac{1}{L_1} + \frac{1}{L_2} + \dots + \frac{1}{Ln}$
    \end{center}
    \item Capacitance(C)
    \begin{center}
        Series: $ \frac{1}{C_{eq}} = \frac{1}{C_1} + \frac{1}{C_2} + \frac{1}{C_n}$\\
        Parallel: $ C_{eq} = C_1 + C_2 + \dots + C_n$
    \end{center}
    \item RMS Value(Any Periodic Wave): $$ V_{rms} = \sqrt{{\frac{1}{T}}\int_0^Tv^2(t)dt}$$
    \item Average Value(Any periodic Wave): $$ V_{avg} =\frac{1}{T}\int_0^Tv(t)dt$$
    \item Complex of AC Circuits:\\
    Rectangular to Polar Conversion: 
    $|Z| = \sqrt{R^2 + X^2},\  \theta = tan^{-1}(\frac{X}{R})$
    \item Polar to Rectangular Conversion: $$ R = |Z|cos\theta, \ X = |Z|sin\theta$$
    \item Admittance Components Explicitly:\\
    If $Z = R + jX$, then $ Y = \frac{1}{Z} = G + jB $
    \begin{center}
        $ G = \frac{R}{R^2 + X^2} \Rightarrow (Conductance) $ \\
        $ B = \frac{-X}{R^2 + X^2} \Rightarrow (Susceptance) $
    \end{center}
\end{itemize}
\section{AC Voltage, Current and Power(R, L, C, R-L, R-C, R-L-C)}
\begin{itemize}
    \item Purely Resistive(R) Circuit:
    \begin{itemize}
        \item Impedance: Z = R\angle 0
        \item Phase: Voltage and Current are in phase($\theta = 0)$
        \item Power Factor: PF = cos(0) = 1
    \end{itemize}
    \item Purely Inductive(L) Circuit:
    \begin{itemize}
        \item Impedance: Z = $ \omega L\angle 90 = j\omega L $
        \item Phase: Current lags voltage by 90 (0 = 90).
        \item Power Factor: PF = cos(90) = 0(lagging)
    \end{itemize}
    \item Purely Capacitive(C) Circuit:
    \begin{itemize}
        \item $Impedance: Z = \frac{1}{\omega C}\angle -90 = -j\frac{1}{\omega C}$
        \item Phase: Current leads voltage by 90($\theta = -90$
        \item Power Factor: PF = cos(-90) = 0(leading)
    \end{itemize}
    \item Series R-L Circuit:
    \begin{itemize}
        \item Impedance: $ Z = R + j\omega L = \sqrt{R^2 + (\omega L)^2}\angle tan^{-1}(\frac{\omega L}{R})$
        \item Current lags voltage by angle $\theta$
    \end{itemize}
    \item Series R-C Circuit:
    \begin{itemize}
        \item Impedance: Z = R -$ j\frac{1}{\omega C} = \sqrt{R^2 +(\frac{1}{\omega C})^2}\angle -tan^{-1}(\frac{1}{\omega CR})$
    \end{itemize}
    
\end{itemize}
\section{AC Series-Parallel Simplification(R, L, C)}
\begin{itemize}
    \item Impedance in Series: $$ Z_{eq} = Z_1 + Z_2 + Z_3+\dots$$
    \item Impedance in Parallel: $\frac{1}{Z_{eq}} = \frac{1}{Z_1} + \frac{1}{Z_2} + \dots or for two component : Z_{eq} = \frac{Z_1\cdot Z_2}{Z_1 + Z_2}$
    \item AC Delta-Wye \ ($\delta-Y) Conversion$ \\
    The conversion formulas are structurally identical to the DC resistance formulas, but you use Complex Impedance(Z) instead of just R:
    \begin{itemize}
        \item Delta ($\Delta$) to Wye (Y):
\begin{align*}
Z_1 &= \frac{Z_bZ_c}{Z_a + Z_b + Z_c},\\
Z_2 &= \frac{Z_aZ_c}{Z_a + Z_b + Z_c},\\
Z_3 &= \frac{Z_aZ_b}{Z_a + Z_b + Z_c}.
\end{align*}
       \begin{align*}
Z_a &= \frac{Z_1Z_2 + Z_2Z_3 + Z_3Z_1}{Z_1},\\
Z_b &= \frac{Z_1Z_2 + Z_2Z_3 + Z_3Z_1}{Z_2},\\
Z_c &= \frac{Z_1Z_2 + Z_2Z_3 + Z_3Z_1}{Z_3}.
\end{align*}
        \item AC Maximum Power Transfer
\begin{itemize}
    \item Reactive part $X_L$ is fixed, and the magnitude $|Z_L|$ is varied.
    \item
    \[
    R_L=\sqrt{R_{Th}^{2}+\left(X_{Th}+X_L\right)^{2}}
    \]
    \item If the circuit is restricted to matching only the magnitudes,
    \[
    |Z_L|=|Z_{Th}|=\sqrt{R_{Th}^{2}+X_{Th}^{2}}
    \]
\end{itemize}
    \end{itemize}
   
\end{itemize}

\end{enumerate}
\end{document}
